\documentclass[../../syllabus_1.tex]{subfiles}
\begin{document}

\chapter{Solides}

\faBookOpen~ Pages 63 à 69.

\section{Caractéristiques des solides}

Un solide est une figure à trois dimensions limitée par une surface fermée. Généralement, il possède des faces, des arêtes et des sommets.

\begin{center}
    \includegraphics{chapitres/solides/tikz/cube_annote.tikz}
\end{center}

En géométrie, l’apex est le nom parfois donné à un sommet distingué des autres. Par exemple, le sommet opposé à la base d'un triangle, d'une pyramide.

Notion de hauteur + dénombrement

\begin{exercise}
    Nomme les solides ci-dessous et dénombre les faces, les arêtes et les sommets.
    
    \begin{tblr}{
            hlines,
            vlines,
            colspec={lX[c]X[c]X[c]X[c]X[c]X[c]},
            % columns={1.5cm, c},
            rows={1cm, m},
            % column{1}={2cm, l},
            vline{1}={1}{white},
            hline{1}={1}{white}
        }
        &
        \includegraphics[width=1cm]{chapitres/solides/tikz/sphere.tikz}
        &
        \includegraphics[width=1cm]{chapitres/solides/tikz/cone.tikz}
        &
        \includegraphics[width=1cm]{chapitres/solides/tikz/tetraedre.tikz}
        &
        \includegraphics[width=1cm]{chapitres/solides/tikz/cylindre.tikz}
        &
        \includegraphics[width=1cm]{chapitres/solides/tikz/parallelepipede_rectangle.tikz}
        &
        \includegraphics[width=1cm]{chapitres/solides/tikz/prisme_droit_hexagone.tikz}\\
        Noms \\
        Faces                    \\
        Arêtes                   \\
        Sommets
    \end{tblr}
\end{exercise}

\section{Classement des solides}

\textbf{Réalisation d'un schéma dans lequel on classe les différents solides.}


Parmi les solides que tu viens de rencontrer certains possèdes des faces planes et d'autre des faces courbes. On utilise ce critère pour différencier les polyèdres et les non-polyèdres. Autrement dit, les non-polyèdres sont des solides qui peuvent rouler, car ils possèdent une face courbe à l'inverse des polyèdres.

Un prisme est un solide géométrique délimité par deux polygones, appelés les bases du prisme, images l'un de l'autre par une translation. Ces bases sont reliées entre elles par des parallélogrammes.

Quand ces parallélogrammes sont des rectangles, on dit que le prisme est droit.


\includegraphics{chapitres/solides/img/solides.png}

Ne pas oublier de parler des polyèdres réguliers.

\section{Représentation de solides}
\subsection{Perspective cavalière}
Lis la théorie page 65 puis réalise l'exercice.
\begin{exercise}
    Voici la face avant d'un cube. Représente le cube en perspective cavalière avec un angle de fuite de \ang{45} et un rapport de réduction de \(\dfrac{2}{3}\).

   \includegraphics{chapitres/solides/tikz/cube_cavaliere.tikz}
   
\end{exercise}
\subsection{Perspective à un point de fuite}
Lis la théorie page 65 puis réalise l'exercice.
\begin{exercise}
    Voici la face avant d'un cube. Représente le cube en perspective à un point de fuite, le point \(F\).

    \includegraphics{chapitres/solides/tikz/cube_fuite.tikz}
\end{exercise}
\section{Développement de solides}
\section{Vues coordonnées}

\section{Positions relatives d'un plan et d'une droite}
\section{Positions relatives de deux droites}
\section{Périmètre, aire et volume}

\end{document}